\chapter{分子结构和分子光谱}

\problem{5.1}{KI 分子原子核的平衡间距为 $0.305\ \mathrm{nm}$，K 原子的第一电离能为 $4.34\ \mathrm{eV}$，I 原子的电子亲和能为 $3.08\ \mathrm{eV}$，试计算中性 K 原子和 I 原子形成 KI 分子时的结合能。}
% layout
\answerarea

\problem{5.2}{K 原子的第一电离能为 $4.34\ \mathrm{eV}$，Cl 原子的电子亲和能为 $3.82\ \mathrm{eV}$，KCl 分子解离为一个 K 原子和一个 Cl 原子需能量为 $4.64\ \mathrm{eV}$，求 KCl 分子的远红外第一条转动光谱线的波数。}
% layout
\answerarea
\newpage

\problem{5.3}{HBr 分子的远红外吸收光谱是一些 $\Delta\bar\nu=16.94\ \mathrm{cm^{-1}}$ 等间隔的光谱线，试求 HBr 分子的转动惯量及原子核的距离。}
% layout
\answerarea

\problem{5.4}{试以 $\mathrm H^{35}\mathrm{Cl}$ 分子为例说明，室温下大部分分子处于振动基态。（设 H 原子与 Cl 原子之间相互作用力常数为 $516\ \mathrm{N\,m^{-1}}$。）}
% layout
\answerarea
\newpage

\problem{5.5}{已知 $\mathrm H^{35}\mathrm{Cl}$ 分子的基频带波数为 $2886.19\ \mathrm{cm^{-1}}$，求：
\begin{enumerate}
\item $\mathrm H^{37}\mathrm{Cl}$ 分子相应的基频谱带的波数；
\item 用波数表示出由这种同位素效应所引起的谱线之间的间距。
\end{enumerate}}
% layout
\answerarea

\problem{5.6}{$^{12}\mathrm C^{16}\mathrm O$ 分子基频带的谱线间隔近似为常数 $3.86\ \mathrm{cm^{-1}}$，谱带中心缺失的谱线波数为 $2170.21\ \mathrm{cm^{-1}}$，试求转动常数 $B$、核间距 $R_0$ 和振动力常数 $k$。}
% layout
\answerarea
\newpage

\problem{5.7}{\begin{enumerate}
\item 假定 $^{23}\mathrm{Na}^{35}\mathrm{Cl}$ 和 $^{23}\mathrm{Na}^{37}\mathrm{Cl}$ 两种同位素分子的核间距相等，一台光谱仪恰能分辨这两种同位素的转动光谱，求这一光谱仪的分辨率 $\Delta\lambda/\lambda$；
\item 若这两种同位素分子的力常数相等，试问这台光谱仪是否也能分辨这两种分子的振动光谱？
\end{enumerate}}
% layout
\answerarea

\problem{5.8}{N$_2$ 分子某激发态的电子组态是
\[
\begin{aligned}
&(\sigma_{\mathrm g}1\mathrm s)^2(\sigma_{\mathrm u}^{*}1\mathrm s)^2
(\sigma_{\mathrm g}2\mathrm s)^2(\sigma_{\mathrm u}^{*}2\mathrm s)^2\\[-2pt]
&\qquad\times(\pi_{\mathrm u}2\mathrm p)^4(\sigma_{\mathrm g}2\mathrm p)^1
(\sigma_{\mathrm u}2\mathrm p)^1.
\end{aligned}
\]
请写出该电子组态按照罗素--桑德尔斯耦合得到可能的分子电子态。}
% layout
\answerarea
\newpage

\problem{5.9}{N$_2$ 分子处在电子基态 $^1\Sigma_{\mathrm g}$ 和激发态 $^1\Sigma_{\mathrm u}$ 的常数如下（单位为 $\mathrm{cm^{-1}}$）：
\begin{center}
\begin{tabular}{lrrr}
\toprule
 & $\bar\nu_0$ & $\gamma\bar\nu_0$ & $B$\\
\midrule
基态 $^1\Sigma_{\mathrm g}$ & 2359.6 & 14.456 & 2.010\\
激发态 $^1\Sigma_{\mathrm u}$ & 751.7 & 4.820 & 1.145\\
\bottomrule
\end{tabular}
\end{center}
这两个电子态之间的跃迁对应的波数为 $103678.9\ \mathrm{cm^{-1}}$。
\begin{enumerate}
\item 计算该电子振转光谱中，$(0,0)$ 带的 R 支和 P 支各自前四条谱线的波数；
\item 说明 $(0,0)$ 带的带头在红端还是在紫端？并求出带头的位置。
\end{enumerate}}
% layout
\answerarea
\newpage

\problem{5.10}{用氦氖激光器波长为 $632.8\ \mathrm{nm}$ 的激光照射 HCl 样品，在散射光中出现了 $532.1\ \mathrm{nm}$ 和 $780.5\ \mathrm{nm}$ 两种波长的谱线。试解释这两条谱线产生的原因，并据此计算 HCl 分子的经典振动频率（不考虑非简谐振动）。}
% layout
\answerarea

\newpage

